\documentclass[UTF8,11pt,fontset=none]{ctexart}
\usepackage{amsmath,amssymb,amsthm,geometry,hyperref,fontspec,xeCJK}
\setmainfont{Times New Roman}\setCJKmainfont{Songti SC}\setCJKsansfont{Heiti SC}
\geometry{a4paper,left=25mm,right=25mm,top=23mm,bottom=24mm}
\theoremstyle{definition}\newtheorem{theorem}{定理}[section]\newtheorem{lemma}[theorem]{引理}\newtheorem{corollary}[theorem]{推论}\newtheorem{proposition}[theorem]{命题}
\theoremstyle{remark}\newtheorem{remark}[theorem]{说明}
\newcommand{\R}{\mathbb R}\newcommand{\1}{\mathbf 1}
\title{中心方体的两种加权 Solyanik 模型\\共同中心多线性精确传递与端点结论}
\author{基于《课题2（8.26最新扩充版）》的核验修订}\date{2026 年 11 月 5 日}
\begin{document}\maketitle
\begin{abstract}本文区分内禀权平均与外层加权光晕，证明共同中心多线性传递、内禀端点和外层集中模量界，并得到不等指数近端点的显式公式。对有限总质量权及任意维乘积权，进一步证明内禀多集合模型的临界发散。\end{abstract}

\section{引言}
中心极大算子的近单位水平集刻画了平均几乎等于一时的外溢质量。加权情形存在两个不同问题：权可以同时进入平均与测量，也可以只进入测量。Hagelstein--Parissis 的加权研究说明这种区分对于端点估计十分重要\cite{HagelsteinParissis2016}；Aldaz 的测度一致弱型理论则提供中心方体的几何工具\cite{Aldaz2020}。
本文围绕共同中心、共同尺度下的多集合输入展开。核心是把各输入的总缺陷与公共交集联系起来，并通过有限稳定见证和同时分割证明反向传递。随后给出逐权一阶项、全权族精确公式和外层集中模量界。进一步用水位优化把等指数传递推广到不等正指数。本文不以局部模型的系数替代完整算子常数，也不声称这些推论具有未经文献比较的首创性。
\section{预备知识与主结果}
令 $Q(x,r)=x+[-r,r]^n$，$d\mu=w\,dx$，其中 $0\le w\in L^1_{\rm loc}(\R^n)$ 且 $w\ne0$。定义
\[
 M_{c,w}f=\sup_{r:\mu(Q)>0}\mu(Q)^{-1}\int_Q|f|d\mu,\quad
 M_cf=\sup_{r>0}|Q|^{-1}\int_Q|f|dx,
\]
\[
 C_{c,w}^{\rm int}(a)=\sup_{0<\mu(E)<\infty}\frac{\mu\{M_{c,w}\1_E>a\}}{\mu(E)},\quad
 C_{c,w}^{\rm out}(a)=\sup_{\substack{0<|E|<\infty\\0<\mu(E)<\infty}}\frac{\mu\{M_c\1_E>a\}}{\mu(E)}.
\]
内禀上确界忽略零质量方体；外层定义保留 $|E|$ 的条件。

令整数 $m\ge2$、$F=\cap_{i=1}^mE_i$，并定义
\[
 \mathcal M^{(m)}_{c,w}(\1_{E_1},\ldots,\1_{E_m})=
 \sup_r\prod_{i=1}^m\frac{\mu(E_i\cap Q)}{\mu(Q)}.
\]
对所有可测 $E_1,\ldots,E_m$ 且 $0<\mu(F)<\infty$，以 $\mu(F)$ 为分母对光晕测量之比取上确界得到 $C^{(m),\rm int}_{c,n,w}$；将乘积中的 $\mu$ 换成 Lebesgue 测度、而光晕和分母仍用 $\mu$，得到 $C^{(m),\rm out}_{c,n,w}$；外层另要求 $0<|F|<\infty$。各 $E_i$ 本身不要求有限测度，所有算子水平均取严格不等式。写
\[
 s_m(\lambda)=m(1-\lambda^{1/m}),\quad\Phi_m(\lambda)=1-s_m(\lambda),\quad
 \lambda_m^*=(1-1/m)^m.
\]
\begin{theorem}\label{transfer}若 $\lambda_m^*<\lambda<1$，则
\[
C^{(m),\rm int}_{c,n,w}(\lambda)=C^{\rm int}_{c,w}(\Phi_m(\lambda)),\qquad
C^{(m),\rm out}_{c,n,w}(\lambda)=C^{\rm out}_{c,w}(\Phi_m(\lambda)).
\]
\end{theorem}
\begin{theorem}\label{endpoint}对每个非零 $w\in L^1_{\rm loc}$，
\[C^{\rm int}_{c,w}(1-t)=1+2^nt+o_w(t)\quad(t\downarrow0).\]
并且对 $0<a<1$，\[\sup_wC^{\rm int}_{c,w}(a)=1+2^n\frac{1-a}{a}.\]
\end{theorem}
\begin{theorem}\label{outer}令
\[\Omega_w(t)=\sup_{\substack{A\subset Q,\ |A|\le t|Q|\\Q\ \mathrm{centered},\ \mu(Q)>0}}\frac{\mu(A)}{\mu(Q)}.\]
若 $\Omega_w(t)<1$，则\[C^{\rm out}_{c,w}(1-t)\le1+2^n\frac{\Omega_w(t)}{1-\Omega_w(t)}.\]
\end{theorem}
\section{共同中心传递的证明}
中心方体的 Besicovitch intersection constant 是 $2^n$，并给出任意局部有限 Radon 测度的测度一致中心弱型 $(1,1)$ 界 \cite[Theorem 1.2 and Corollary 3.5]{Aldaz2020}。
\begin{lemma}\label{deficit}若 $0\le a_i\le1$、$\prod_ia_i>\lambda$，则
\[\sum_i(1-a_i)<s_m(\lambda).\]
该常数最优，且 $\Phi_m(\lambda)>0$ 当且仅当 $\lambda>\lambda_m^*$。
\end{lemma}
\begin{proof}AM--GM 给出 $m^{-1}\sum_ia_i\ge(\prod_ia_i)^{1/m}>\lambda^{1/m}$。移项即得；取全部 $a_i$ 相等即知最优性。\end{proof}
\begin{lemma}\label{split}设 $\nu$ 为 Lebesgue 测度或 $w\,dx$，$F\subset D$、$\nu(D)<\infty$，且 $Q_1,\ldots,Q_N\subset D$。则 $D\setminus F$ 可分成不交的 $A_1,\ldots,A_m$，使
\[\nu(A_i\cap Q_j)=m^{-1}\nu((D\setminus F)\cap Q_j).\]
\end{lemma}
\begin{proof}由有限方体族生成 Boolean 原子。在每个正质量原子上沿第一坐标按连续分布函数等分；逐原子求和。\end{proof}
\begin{proof}[定理 \ref{transfer} 的证明]
若同一方体上的输入比例为 $a_i$，引理 \ref{deficit} 和 $Q\setminus F\subset\cup_i(Q\setminus E_i)$ 给出
\[\nu(F\cap Q)/\nu(Q)>1-s_m(\lambda)=\Phi_m(\lambda),\]
其中内禀取 $\nu=\mu$、外层取 $\nu=dx$，故得到上界。

反向构造是第 \ref{sec:weighted-exponents} 节定理 \ref{thm:weighted-transfer} 在 $a_i=1$ 时的情形。该处详细处理有限参考见证的中心扰动、Boolean 原子的同时等分，以及 $D$ 外的分割延拓；特别是构造始终满足 $\cap_iE_i=F$。该证明不使用本定理，因而没有循环依赖。\end{proof}
\begin{corollary}在 $\lambda_m^*<\lambda<1$，
\[\sup_wC^{(m),\rm int}_{c,n,w}(\lambda)=1+2^n\frac{s_m(\lambda)}{1-s_m(\lambda)}.\]
\end{corollary}
\section{内禀端点与外层集中模量}
对紧 $K\subset\{M_{c,w}\1_E>a\}\setminus E$，令 $B=(1-a)/a$ 和 $d\nu_c=\1_Ed\mu+c\1_Kd\mu$。每个严格 $\mu$-见证给出 $M_{c,\nu_c}\1_E>1/(1+cB)$ 于 $K$。弱型界遂给
\[\mu(K)\le2^n(B+c^{-1})\mu(E).\]
令 $c\to\infty$ 并用内正则性，得到
\[C^{\rm int}_{c,w}(a)\le1+2^n(1-a)/a.\tag{1}\]
\subsection{周期穿孔与固定权下界}
固定 $N\ge2$ 与 $0<\theta<1/2$，令
\[
 h=\frac{\theta}{N(1+2\theta)},\qquad r=\frac1{N(1+2\theta)},\qquad D=[0,1]^n.
\]
在每个 $k/N$，$k\in\{1,\ldots,N-1\}^n$，挖去半边长为 $h$ 的开孔 $P_k$，记剩余集为 $E_{N,\theta}$。任意孔点 $z$ 的方体 $Q(z,r)$ 位于 $D$，包含整个本孔，且只可能在边界接触其他孔，因为 $r\ge2h$、$r+2h=1/N$。故该见证的缺失比例为 $(h/r)^n=\theta^n$。孔外点除有限超平面边界外都是剩余集的内点。因此当 $\theta^n<t$ 时，$D$ 模零包含于 $\{M_c\1_{E_{N,\theta}}>1-t\}$，且
\[
 |E_{N,\theta}|=1-(N-1)^n(2h)^n.
\]
对内禀固定权，取正有限 Lebesgue 点 $x_0$。对所有可测 $A\subset D$，有统一估计
\[
 \left|\rho^{-n}\mu(x_0+\rho A)-w(x_0)|A|\right|
 \le\int_D|w(x_0+\rho y)-w(x_0)|\,dy=:\varepsilon_\rho\longrightarrow0.
\]
固定 $N,\theta$ 后，全部孔点的见证体积都等于 $q=(2r)^n>0$。其缩放后的内禀平均至少为
\[
 \frac{w(x_0)(1-\theta^n)q-\varepsilon_\rho}
 {w(x_0)q+\varepsilon_\rho}>1-t
\]
当 $\rho$ 充分小时成立；孔外剩余集的内部平均为一，边界对 $\mu\ll dx$ 为零集。因此先令 $\rho\downarrow0$、再令 $N\to\infty$、最后令 $\theta\uparrow t^{1/n}$，得到对 $0<t<2^{-n}$ 的下包络
\[
 C^{\rm int}_{c,w}(1-t)\ge L_n(t):=
 \left[1-\left(\frac{2t^{1/n}}{1+2t^{1/n}}\right)^n\right]^{-1}.
\]
由于 $L_n(t)=1+2^nt+O_n(t^{1+1/n})$，与 (1) 合并即得定理 \ref{endpoint} 的第一式。这里每个 $t$ 允许重新选择缩放参数，不要求权的收敛速度一致。特别地，$n\ge2$ 时仍有可靠的二阶夹逼
\[
 -n2^{n+1}\le\liminf_{t\downarrow0}\frac{C^{\rm int}_{c,w}(1-t)-1-2^nt}{t^{1+1/n}}
 \le\limsup_{t\downarrow0}\frac{C^{\rm int}_{c,w}(1-t)-1-2^nt}{t^{1+1/n}}\le0.
\]
这不等于二阶极限存在。
\subsection{全权族的匹配下界}
为证 (1) 对权族的下界，取 $B_0=[-\delta,\delta]^n$、$E=B_0$、$B_v=v+B_0$，其中 $v\in\{-1,1\}^n$、$\delta<1/4$。令权在中心块质量为 $1$，每个卫星块质量为 $q=(1-\eta)(1-a)/a$，其他地方为零。每个卫星点的中心方体 $Q(x,1+2\delta)$ 恰捕获中心与该卫星，平均 $1/(1+q)>a$，故光晕商至少 $1+2^nq$。令 $\eta\downarrow0$ 即得定理 \ref{endpoint} 的第二式。

对外层模型，取 $K\Subset\{M_c\1_E>1-t\}\setminus E$，记 $\Omega=\Omega_w(t)$，仍令 $d\nu_c=\1_Ed\mu+c\1_Kd\mu$。Lebesgue 严格见证给出
\[\mu(E\cap Q)\ge(1-\Omega)\mu(Q),\qquad\mu(K\cap Q)\le\Omega\mu(Q).\]
故 $M_{c,\nu_c}\1_E\ge[1+c\Omega/(1-\Omega)]^{-1}$ 于 $K$。对任意略小的阈值应用弱型界再取极限。弱型界、$c\to\infty$ 和内正则性给出定理 \ref{outer}。若 $\Omega_w(t)\le A_0t^\sigma\le1/2$，右端减一至多为 $2^{n+1}A_0t^\sigma$。

% 可由 % 可由 % 可由 \input{extension.tex} 纳入主稿；要求 amsmath, amssymb, amsthm。
\section{不等指数的共同中心传递}\label{sec:weighted-exponents}

本节把等指数乘积的缺陷坐标推广到任意正指数。该推广仍要求各因子在同一个中心方体上取平均；分别取上确界的多线性模型不适用本节结论。

令 $m\ge2$、$a_1,\ldots,a_m>0$，并在 $0\le s\le1$ 定义
\[
 P_{\boldsymbol a}(s)=
 \max\left\{\prod_{i=1}^m(1-d_i)^{a_i}:
 0\le d_i<1,\ \sum_{i=1}^m d_i=s\right\}.
 \tag{W.1}
\]
函数 $P_{\boldsymbol a}$ 连续、严格递减，且 $P_{\boldsymbol a}(0)=1$。因此当
\[
 P_{\boldsymbol a}(1)<\lambda<1
 \tag{W.2}
\]
时，存在唯一的 $s_{\boldsymbol a}(\lambda)\in(0,1)$ 满足
\[
 P_{\boldsymbol a}(s_{\boldsymbol a}(\lambda))=\lambda.
 \tag{W.3}
\]

\begin{lemma}[水位分配]\label{lem:water-filling}
若 $0<s<1$，存在唯一的 $\rho>0$ 使
\[
 \sum_{i=1}^m\left(1-\frac{a_i}{\rho}\right)_+=s.
 \tag{W.4}
\]
式 \textup{(W.1)} 的唯一最大点为
\[
 d_i^*(s)=\left(1-\frac{a_i}{\rho}\right)_+.
 \tag{W.5}
\]
\end{lemma}

\begin{proof}
在单纯形的相对内部，目标的对数为严格凹函数
\(\sum_i a_i\log(1-d_i)\)。KKT 条件给出：活跃坐标满足
\(a_i/(1-d_i)=\rho\)，非活跃坐标满足 $a_i\ge\rho$。于是得到
\textup{(W.4)}--\textup{(W.5)}。左边作为 $\rho$ 的函数连续递增，从 $0$ 取到大于 $1$ 的值，因而 $\rho$ 唯一。严格递减性和连续性也随之得到：增加总缺陷后，把一个正缺陷坐标减小即可严格增大乘积。
\end{proof}

设 $\mu$ 是非原子、局部有限的 Radon 测度，零质量方体一律略去，
\[
 \mathcal M_{c,\mu}^{\boldsymbol a}(f_1,\ldots,f_m)(x)
 =\sup_{r>0}\prod_{i=1}^m
 \left(\frac{1}{\mu(Q(x,r))}\int_{Q(x,r)}|f_i|\,d\mu\right)^{a_i}.
\]
定义共同中心常数
\[
 C_{c,\mu}^{\boldsymbol a}(\lambda)=
 \sup_{\substack{E_1,\ldots,E_m\ \mathrm{measurable}\\
 F=\cap_iE_i,\ 0<\mu(F)<\infty}}
 \frac{\mu\{\mathcal M_{c,\mu}^{\boldsymbol a}
 (\mathbf1_{E_1},\ldots,\mathbf1_{E_m})>\lambda\}}{\mu(F)}.
\]

\begin{theorem}[不等指数共同中心精确传递]\label{thm:weighted-transfer}
在 \textup{(W.2)} 的区间内，
\[
 C_{c,\mu}^{\boldsymbol a}(\lambda)=
 C_{c,\mu}^{\mathrm{int}}\bigl(1-s_{\boldsymbol a}(\lambda)\bigr).
 \tag{W.6}
\]
若 $d\mu=w\,dx$，上式适用于任意非零非负局部可积权。对外层模型，把左侧平均改为 Lebesgue 平均、光晕测量改为 $w\,dx$，并要求 $0<|F|<\infty$，仍有同一等式。
\end{theorem}

\begin{proof}
写 $s=s_{\boldsymbol a}(\lambda)$。若在共同见证方体 $Q$ 上
\[
 p_i=\frac{\mu(E_i\cap Q)}{\mu(Q)},\qquad d_i=1-p_i,
\]
且 $\prod_i p_i^{a_i}>\lambda$，若总缺陷至少为 $s$，可逐坐标减小缺陷到总和恰为 $s$；乘积只能增大，且至多为 $P_{\boldsymbol a}(s)=\lambda$，矛盾。因此定义 \textup{(W.1)} 给出
\(\sum_i d_i<s\)。又 $Q\setminus F\subset\bigcup_i(Q\setminus E_i)$，故
\[
 \frac{\mu(F\cap Q)}{\mu(Q)}
 \ge1-\sum_i d_i>1-s.
\]
这给出 \textup{(W.6)} 的“$\le$”方向。

反之，令 $K$ 为严格线性光晕
\(\{M_{c,\mu}\mathbf1_F>1-s\}$ 的紧子集。我们使用如下有限稳定事实：可选有限个参考方体以及正余量，使 $K$ 被相应的小中心邻域覆盖；每个邻域中的中心方体仍有 $F$-缺陷严格小于 $s$；并且参考方体及这些小中心扰动的对称差的 $\mu$-质量小到小于预先指定的余量。先把严格见证半径略微增大，避开至多可数个具有正质量边界的半径，使每个参考方体边界的 $\mu$-质量为零。严格性在充分小的增大下保持。此时这一事实由严格不等式、局部有限性、从上连续性和
\(\mu(Q(x,r)\mathbin\triangle Q(y,r))\to0\) 得到。

取包含全部微扰方体的有界盒 $D$。把有限参考方体所生成的 Boolean 原子上的 $D\setminus F$ 按下列比例切分：
\[
 \frac{d_1^*(s)}s,\ldots,\frac{d_m^*(s)}s.
\]
非原子性保证每个原子可作这种有限分割；在 $D$ 外把 $F^c$ 任意补入这些部分，得到 $F^c=A_1\mathbin{\dot\cup}\cdots\mathbin{\dot\cup}A_m$，并置 $E_i=\mathbb R^n\setminus A_i$。对未扰动的参考方体，
\[
 \frac{\mu(Q\setminus E_i)}{\mu(Q)}=
 \frac{d_i^*(s)}s\frac{\mu(Q\setminus F)}{\mu(Q)}<d_i^*(s)\quad(d_i^*(s)>0).
 \tag{W.7}
\]
有限稳定性中的余量使 \textup{(W.7)} 在所覆盖的全部中心扰动方体中仍严格成立。这里把 $D$ 取为包含所有微扰方体的一个有界盒；在该盒内按参考方体的 Boolean 原子分割。若 $d_i^*(s)=0$，令 $A_i\cap D=\varnothing$，故该坐标没有误差。于是每个相关方体满足
\[
 \prod_i\left(\frac{\mu(E_i\cap Q)}{\mu(Q)}\right)^{a_i}
 >\prod_i(1-d_i^*(s))^{a_i}=\lambda.
\]
因此 $K$ 包含于所构造的多线性光晕。由 Radon 内正则性令 $K$ 穷尽线性光晕，得到反向不等式。

外层版本完全相同：对 Lebesgue 测度作有限 Boolean 原子切分，再以 $w\,dx$ 的 Radon 内正则性穷尽光晕。
\end{proof}

\begin{corollary}[等指数特例]
当 $a_1=\cdots=a_m=1$ 时，
\[
 P_{\boldsymbol a}(s)=(1-s/m)^m,\qquad
 s_{\boldsymbol a}(\lambda)=m(1-\lambda^{1/m}),
\]
从而恢复等指数的共同中心缺陷传递。
\end{corollary}

\begin{remark}[外层临界发散]
若 $m\ge2$，则对每个非零局部可积权 $w$ 和
\(0<\lambda\le P_{\boldsymbol a}(1)\)，外层常数为无穷。证明取沿第一坐标的周期分割，其各部分的比例为 $d_i^*(1)$；把一个有正 $w$-质量的有限集 $F$ 放进某个正比例部分，并令 $E_i=(\mathbb R^n\setminus A_i)\cup F$。每个充分大的整数边长中心方体给出严格大于 $P_{\boldsymbol a}(1)$ 的共同见证，故光晕为全空间。有限总质量情形再令 $w(F)\downarrow0$。该论证改变内禀平均分母，不能推出内禀模型的临界发散。
\end{remark}
 纳入主稿；要求 amsmath, amssymb, amsthm。
\section{不等指数的共同中心传递}\label{sec:weighted-exponents}

本节把等指数乘积的缺陷坐标推广到任意正指数。该推广仍要求各因子在同一个中心方体上取平均；分别取上确界的多线性模型不适用本节结论。

令 $m\ge2$、$a_1,\ldots,a_m>0$，并在 $0\le s\le1$ 定义
\[
 P_{\boldsymbol a}(s)=
 \max\left\{\prod_{i=1}^m(1-d_i)^{a_i}:
 0\le d_i<1,\ \sum_{i=1}^m d_i=s\right\}.
 \tag{W.1}
\]
函数 $P_{\boldsymbol a}$ 连续、严格递减，且 $P_{\boldsymbol a}(0)=1$。因此当
\[
 P_{\boldsymbol a}(1)<\lambda<1
 \tag{W.2}
\]
时，存在唯一的 $s_{\boldsymbol a}(\lambda)\in(0,1)$ 满足
\[
 P_{\boldsymbol a}(s_{\boldsymbol a}(\lambda))=\lambda.
 \tag{W.3}
\]

\begin{lemma}[水位分配]\label{lem:water-filling}
若 $0<s<1$，存在唯一的 $\rho>0$ 使
\[
 \sum_{i=1}^m\left(1-\frac{a_i}{\rho}\right)_+=s.
 \tag{W.4}
\]
式 \textup{(W.1)} 的唯一最大点为
\[
 d_i^*(s)=\left(1-\frac{a_i}{\rho}\right)_+.
 \tag{W.5}
\]
\end{lemma}

\begin{proof}
在单纯形的相对内部，目标的对数为严格凹函数
\(\sum_i a_i\log(1-d_i)\)。KKT 条件给出：活跃坐标满足
\(a_i/(1-d_i)=\rho\)，非活跃坐标满足 $a_i\ge\rho$。于是得到
\textup{(W.4)}--\textup{(W.5)}。左边作为 $\rho$ 的函数连续递增，从 $0$ 取到大于 $1$ 的值，因而 $\rho$ 唯一。严格递减性和连续性也随之得到：增加总缺陷后，把一个正缺陷坐标减小即可严格增大乘积。
\end{proof}

设 $\mu$ 是非原子、局部有限的 Radon 测度，零质量方体一律略去，
\[
 \mathcal M_{c,\mu}^{\boldsymbol a}(f_1,\ldots,f_m)(x)
 =\sup_{r>0}\prod_{i=1}^m
 \left(\frac{1}{\mu(Q(x,r))}\int_{Q(x,r)}|f_i|\,d\mu\right)^{a_i}.
\]
定义共同中心常数
\[
 C_{c,\mu}^{\boldsymbol a}(\lambda)=
 \sup_{\substack{E_1,\ldots,E_m\ \mathrm{measurable}\\
 F=\cap_iE_i,\ 0<\mu(F)<\infty}}
 \frac{\mu\{\mathcal M_{c,\mu}^{\boldsymbol a}
 (\mathbf1_{E_1},\ldots,\mathbf1_{E_m})>\lambda\}}{\mu(F)}.
\]

\begin{theorem}[不等指数共同中心精确传递]\label{thm:weighted-transfer}
在 \textup{(W.2)} 的区间内，
\[
 C_{c,\mu}^{\boldsymbol a}(\lambda)=
 C_{c,\mu}^{\mathrm{int}}\bigl(1-s_{\boldsymbol a}(\lambda)\bigr).
 \tag{W.6}
\]
若 $d\mu=w\,dx$，上式适用于任意非零非负局部可积权。对外层模型，把左侧平均改为 Lebesgue 平均、光晕测量改为 $w\,dx$，并要求 $0<|F|<\infty$，仍有同一等式。
\end{theorem}

\begin{proof}
写 $s=s_{\boldsymbol a}(\lambda)$。若在共同见证方体 $Q$ 上
\[
 p_i=\frac{\mu(E_i\cap Q)}{\mu(Q)},\qquad d_i=1-p_i,
\]
且 $\prod_i p_i^{a_i}>\lambda$，若总缺陷至少为 $s$，可逐坐标减小缺陷到总和恰为 $s$；乘积只能增大，且至多为 $P_{\boldsymbol a}(s)=\lambda$，矛盾。因此定义 \textup{(W.1)} 给出
\(\sum_i d_i<s\)。又 $Q\setminus F\subset\bigcup_i(Q\setminus E_i)$，故
\[
 \frac{\mu(F\cap Q)}{\mu(Q)}
 \ge1-\sum_i d_i>1-s.
\]
这给出 \textup{(W.6)} 的“$\le$”方向。

反之，令 $K$ 为严格线性光晕
\(\{M_{c,\mu}\mathbf1_F>1-s\}$ 的紧子集。我们使用如下有限稳定事实：可选有限个参考方体以及正余量，使 $K$ 被相应的小中心邻域覆盖；每个邻域中的中心方体仍有 $F$-缺陷严格小于 $s$；并且参考方体及这些小中心扰动的对称差的 $\mu$-质量小到小于预先指定的余量。先把严格见证半径略微增大，避开至多可数个具有正质量边界的半径，使每个参考方体边界的 $\mu$-质量为零。严格性在充分小的增大下保持。此时这一事实由严格不等式、局部有限性、从上连续性和
\(\mu(Q(x,r)\mathbin\triangle Q(y,r))\to0\) 得到。

取包含全部微扰方体的有界盒 $D$。把有限参考方体所生成的 Boolean 原子上的 $D\setminus F$ 按下列比例切分：
\[
 \frac{d_1^*(s)}s,\ldots,\frac{d_m^*(s)}s.
\]
非原子性保证每个原子可作这种有限分割；在 $D$ 外把 $F^c$ 任意补入这些部分，得到 $F^c=A_1\mathbin{\dot\cup}\cdots\mathbin{\dot\cup}A_m$，并置 $E_i=\mathbb R^n\setminus A_i$。对未扰动的参考方体，
\[
 \frac{\mu(Q\setminus E_i)}{\mu(Q)}=
 \frac{d_i^*(s)}s\frac{\mu(Q\setminus F)}{\mu(Q)}<d_i^*(s)\quad(d_i^*(s)>0).
 \tag{W.7}
\]
有限稳定性中的余量使 \textup{(W.7)} 在所覆盖的全部中心扰动方体中仍严格成立。这里把 $D$ 取为包含所有微扰方体的一个有界盒；在该盒内按参考方体的 Boolean 原子分割。若 $d_i^*(s)=0$，令 $A_i\cap D=\varnothing$，故该坐标没有误差。于是每个相关方体满足
\[
 \prod_i\left(\frac{\mu(E_i\cap Q)}{\mu(Q)}\right)^{a_i}
 >\prod_i(1-d_i^*(s))^{a_i}=\lambda.
\]
因此 $K$ 包含于所构造的多线性光晕。由 Radon 内正则性令 $K$ 穷尽线性光晕，得到反向不等式。

外层版本完全相同：对 Lebesgue 测度作有限 Boolean 原子切分，再以 $w\,dx$ 的 Radon 内正则性穷尽光晕。
\end{proof}

\begin{corollary}[等指数特例]
当 $a_1=\cdots=a_m=1$ 时，
\[
 P_{\boldsymbol a}(s)=(1-s/m)^m,\qquad
 s_{\boldsymbol a}(\lambda)=m(1-\lambda^{1/m}),
\]
从而恢复等指数的共同中心缺陷传递。
\end{corollary}

\begin{remark}[外层临界发散]
若 $m\ge2$，则对每个非零局部可积权 $w$ 和
\(0<\lambda\le P_{\boldsymbol a}(1)\)，外层常数为无穷。证明取沿第一坐标的周期分割，其各部分的比例为 $d_i^*(1)$；把一个有正 $w$-质量的有限集 $F$ 放进某个正比例部分，并令 $E_i=(\mathbb R^n\setminus A_i)\cup F$。每个充分大的整数边长中心方体给出严格大于 $P_{\boldsymbol a}(1)$ 的共同见证，故光晕为全空间。有限总质量情形再令 $w(F)\downarrow0$。该论证改变内禀平均分母，不能推出内禀模型的临界发散。
\end{remark}
 纳入主稿；要求 amsmath, amssymb, amsthm。
\section{不等指数的共同中心传递}\label{sec:weighted-exponents}

本节把等指数乘积的缺陷坐标推广到任意正指数。该推广仍要求各因子在同一个中心方体上取平均；分别取上确界的多线性模型不适用本节结论。

令 $m\ge2$、$a_1,\ldots,a_m>0$，并在 $0\le s\le1$ 定义
\[
 P_{\boldsymbol a}(s)=
 \max\left\{\prod_{i=1}^m(1-d_i)^{a_i}:
 0\le d_i<1,\ \sum_{i=1}^m d_i=s\right\}.
 \tag{W.1}
\]
函数 $P_{\boldsymbol a}$ 连续、严格递减，且 $P_{\boldsymbol a}(0)=1$。因此当
\[
 P_{\boldsymbol a}(1)<\lambda<1
 \tag{W.2}
\]
时，存在唯一的 $s_{\boldsymbol a}(\lambda)\in(0,1)$ 满足
\[
 P_{\boldsymbol a}(s_{\boldsymbol a}(\lambda))=\lambda.
 \tag{W.3}
\]

\begin{lemma}[水位分配]\label{lem:water-filling}
若 $0<s<1$，存在唯一的 $\rho>0$ 使
\[
 \sum_{i=1}^m\left(1-\frac{a_i}{\rho}\right)_+=s.
 \tag{W.4}
\]
式 \textup{(W.1)} 的唯一最大点为
\[
 d_i^*(s)=\left(1-\frac{a_i}{\rho}\right)_+.
 \tag{W.5}
\]
\end{lemma}

\begin{proof}
在单纯形的相对内部，目标的对数为严格凹函数
\(\sum_i a_i\log(1-d_i)\)。KKT 条件给出：活跃坐标满足
\(a_i/(1-d_i)=\rho\)，非活跃坐标满足 $a_i\ge\rho$。于是得到
\textup{(W.4)}--\textup{(W.5)}。左边作为 $\rho$ 的函数连续递增，从 $0$ 取到大于 $1$ 的值，因而 $\rho$ 唯一。严格递减性和连续性也随之得到：增加总缺陷后，把一个正缺陷坐标减小即可严格增大乘积。
\end{proof}

设 $\mu$ 是非原子、局部有限的 Radon 测度，零质量方体一律略去，
\[
 \mathcal M_{c,\mu}^{\boldsymbol a}(f_1,\ldots,f_m)(x)
 =\sup_{r>0}\prod_{i=1}^m
 \left(\frac{1}{\mu(Q(x,r))}\int_{Q(x,r)}|f_i|\,d\mu\right)^{a_i}.
\]
定义共同中心常数
\[
 C_{c,\mu}^{\boldsymbol a}(\lambda)=
 \sup_{\substack{E_1,\ldots,E_m\ \mathrm{measurable}\\
 F=\cap_iE_i,\ 0<\mu(F)<\infty}}
 \frac{\mu\{\mathcal M_{c,\mu}^{\boldsymbol a}
 (\mathbf1_{E_1},\ldots,\mathbf1_{E_m})>\lambda\}}{\mu(F)}.
\]

\begin{theorem}[不等指数共同中心精确传递]\label{thm:weighted-transfer}
在 \textup{(W.2)} 的区间内，
\[
 C_{c,\mu}^{\boldsymbol a}(\lambda)=
 C_{c,\mu}^{\mathrm{int}}\bigl(1-s_{\boldsymbol a}(\lambda)\bigr).
 \tag{W.6}
\]
若 $d\mu=w\,dx$，上式适用于任意非零非负局部可积权。对外层模型，把左侧平均改为 Lebesgue 平均、光晕测量改为 $w\,dx$，并要求 $0<|F|<\infty$，仍有同一等式。
\end{theorem}

\begin{proof}
写 $s=s_{\boldsymbol a}(\lambda)$。若在共同见证方体 $Q$ 上
\[
 p_i=\frac{\mu(E_i\cap Q)}{\mu(Q)},\qquad d_i=1-p_i,
\]
且 $\prod_i p_i^{a_i}>\lambda$，若总缺陷至少为 $s$，可逐坐标减小缺陷到总和恰为 $s$；乘积只能增大，且至多为 $P_{\boldsymbol a}(s)=\lambda$，矛盾。因此定义 \textup{(W.1)} 给出
\(\sum_i d_i<s\)。又 $Q\setminus F\subset\bigcup_i(Q\setminus E_i)$，故
\[
 \frac{\mu(F\cap Q)}{\mu(Q)}
 \ge1-\sum_i d_i>1-s.
\]
这给出 \textup{(W.6)} 的“$\le$”方向。

反之，令 $K$ 为严格线性光晕
\(\{M_{c,\mu}\mathbf1_F>1-s\}$ 的紧子集。我们使用如下有限稳定事实：可选有限个参考方体以及正余量，使 $K$ 被相应的小中心邻域覆盖；每个邻域中的中心方体仍有 $F$-缺陷严格小于 $s$；并且参考方体及这些小中心扰动的对称差的 $\mu$-质量小到小于预先指定的余量。先把严格见证半径略微增大，避开至多可数个具有正质量边界的半径，使每个参考方体边界的 $\mu$-质量为零。严格性在充分小的增大下保持。此时这一事实由严格不等式、局部有限性、从上连续性和
\(\mu(Q(x,r)\mathbin\triangle Q(y,r))\to0\) 得到。

取包含全部微扰方体的有界盒 $D$。把有限参考方体所生成的 Boolean 原子上的 $D\setminus F$ 按下列比例切分：
\[
 \frac{d_1^*(s)}s,\ldots,\frac{d_m^*(s)}s.
\]
非原子性保证每个原子可作这种有限分割；在 $D$ 外把 $F^c$ 任意补入这些部分，得到 $F^c=A_1\mathbin{\dot\cup}\cdots\mathbin{\dot\cup}A_m$，并置 $E_i=\mathbb R^n\setminus A_i$。对未扰动的参考方体，
\[
 \frac{\mu(Q\setminus E_i)}{\mu(Q)}=
 \frac{d_i^*(s)}s\frac{\mu(Q\setminus F)}{\mu(Q)}<d_i^*(s)\quad(d_i^*(s)>0).
 \tag{W.7}
\]
有限稳定性中的余量使 \textup{(W.7)} 在所覆盖的全部中心扰动方体中仍严格成立。这里把 $D$ 取为包含所有微扰方体的一个有界盒；在该盒内按参考方体的 Boolean 原子分割。若 $d_i^*(s)=0$，令 $A_i\cap D=\varnothing$，故该坐标没有误差。于是每个相关方体满足
\[
 \prod_i\left(\frac{\mu(E_i\cap Q)}{\mu(Q)}\right)^{a_i}
 >\prod_i(1-d_i^*(s))^{a_i}=\lambda.
\]
因此 $K$ 包含于所构造的多线性光晕。由 Radon 内正则性令 $K$ 穷尽线性光晕，得到反向不等式。

外层版本完全相同：对 Lebesgue 测度作有限 Boolean 原子切分，再以 $w\,dx$ 的 Radon 内正则性穷尽光晕。
\end{proof}

\begin{corollary}[等指数特例]
当 $a_1=\cdots=a_m=1$ 时，
\[
 P_{\boldsymbol a}(s)=(1-s/m)^m,\qquad
 s_{\boldsymbol a}(\lambda)=m(1-\lambda^{1/m}),
\]
从而恢复等指数的共同中心缺陷传递。
\end{corollary}

\begin{remark}[外层临界发散]
若 $m\ge2$，则对每个非零局部可积权 $w$ 和
\(0<\lambda\le P_{\boldsymbol a}(1)\)，外层常数为无穷。证明取沿第一坐标的周期分割，其各部分的比例为 $d_i^*(1)$；把一个有正 $w$-质量的有限集 $F$ 放进某个正比例部分，并令 $E_i=(\mathbb R^n\setminus A_i)\cup F$。每个充分大的整数边长中心方体给出严格大于 $P_{\boldsymbol a}(1)$ 的共同见证，故光晕为全空间。有限总质量情形再令 $w(F)\downarrow0$。该论证改变内禀平均分母，不能推出内禀模型的临界发散。
\end{remark}

% Input after extension.tex; requires amsmath and amsthm.
\subsection{不等指数的近端点分段}\label{subsec:weighted-endpoint}

设 $a_*=\min_i a_i$，令 $k=\#\{i:a_i=a_*\}$。若并非所有
指数都等于 $a_*$，记
\[
 a_+=\min\{a_i:a_i>a_*\},\qquad
 \widehat s_*=k\left(1-\frac{a_*}{a_+}\right),\qquad
 s_0=\min\{1,\widehat s_*\}.
\]
若所有指数相等，则约定 $s_0=1$。以下断言在 $0\le s<s_0$ 的范围内成立；
仅当 $\widehat s_*<1$ 时，允许在 $s=\widehat s_*$ 取等号。情形
$\widehat s_*=1$ 只按左极限使用，因而不把临界传递端点纳入下述结论。

\begin{theorem}[第一水位区间]\label{thm:first-water-level}
在上述范围内，水位优化函数满足
\[
 P_{\boldsymbol a}(s)=\left(1-\frac{s}{k}\right)^{ka_*}.
 \tag{E.1}
\]
令 $\lambda_0=(1-s_0/k)^{ka_*}$（若 $s_0=1$ 且 $k=1$，取左极限
\(\lambda_0=0\)）。在 $s_0<1$ 时对 $\lambda_0\le\lambda<1$，在
$s_0=1$ 时对 $\lambda_0<\lambda<1$，其反函数为
\[
 s_{\boldsymbol a}(\lambda)
 =k\left(1-\lambda^{1/(ka_*)}\right).
 \tag{E.2}
\]
若 $d\mu=w\,dx$，则当 $\lambda$ 同时满足上述第一水位阈值条件以及
$P_{\boldsymbol a}(1)<\lambda<1$ 时，定理 \ref{thm:weighted-transfer} 给出
\[
 C^{\boldsymbol a}_{c,\mu}(\lambda)
 =C^{\rm int}_{c,\mu}\!\left(1-
 k(1-\lambda^{1/(ka_*)})\right).
 \tag{E.3}
\]
对于每个固定的非零非负 $w\in L^1_{\rm loc}$，当
$\lambda=1-\varepsilon\uparrow1$ 时，
\[
 C^{\boldsymbol a}_{c,w}(1-\varepsilon)
=1+\frac{2^n}{a_*}\varepsilon+o_{\boldsymbol a,w}(\varepsilon).
 \tag{E.4}
\]
并且当 $\lambda$ 满足同一双重条件时，
\[
 \sup_w C^{\boldsymbol a}_{c,w}(\lambda)
 =1+2^n\frac{k(1-\lambda^{1/(ka_*)})}
 {1-k(1-\lambda^{1/(ka_*)})}.
 \tag{E.5}
\]
\end{theorem}

\begin{proof}
先设存在较大指数。在 $s$ 的第一水位区间，令
\[ \rho=\frac{ka_*}{k-s}. \]
由 $s<s_0\le\widehat s_*$，有 $\rho<a_+$。水位引理 \ref{lem:water-filling}
因此表明只有指数为 $a_*$ 的 $k$ 个坐标活跃，且每个活跃缺陷为
$s/k$，其余缺陷为零。若所有指数均为 $a_*$，水位引理直接给出全部
$k=m$ 个坐标活跃及同一缺陷分配。这证明 \textup{(E.1)}。当
$\widehat s_*<1$ 且 $s=\widehat s_*$ 时，新坐标刚好以零缺陷加入，故
公式在该非临界端点仍成立。
\textup{(E.2)} 直接由
\textup{(E.1)} 反解；\textup{(E.3)} 是定理
\ref{thm:weighted-transfer} 的代入。

取 $\lambda=1-\varepsilon$。二项式展开给出
\[
 k\left(1-(1-\varepsilon)^{1/(ka_*)}\right)
 =\frac{\varepsilon}{a_*}+O_{\boldsymbol a}(\varepsilon^2).
\]
再使用固定权线性端点式
$C^{\rm int}_{c,w}(1-t)=1+2^nt+o_w(t)$，即得
\textup{(E.4)}。最后，将内禀全权族公式
$\sup_wC^{\rm int}_{c,w}(1-t)=1+2^nt/(1-t)$ 代入
\textup{(E.2)}，得到 \textup{(E.5)}。
\end{proof}

\begin{remark}
当 $\widehat s_*<1$ 且 $s>\widehat s_*$ 时，至少一个较大指数的坐标进入水位，故
\textup{(E.1)}--\textup{(E.2)} 不再适用；必须使用完整的分段水位
公式。这里的结论不涉及内禀模型在临界以下的发散，也不把数值验证
当作传递定理的证明。
\end{remark}

\section{内禀临界的可解情形}
% 可在 % 可由 % 可由 \input{extension.tex} 纳入主稿；要求 amsmath, amssymb, amsthm。
\section{不等指数的共同中心传递}\label{sec:weighted-exponents}

本节把等指数乘积的缺陷坐标推广到任意正指数。该推广仍要求各因子在同一个中心方体上取平均；分别取上确界的多线性模型不适用本节结论。

令 $m\ge2$、$a_1,\ldots,a_m>0$，并在 $0\le s\le1$ 定义
\[
 P_{\boldsymbol a}(s)=
 \max\left\{\prod_{i=1}^m(1-d_i)^{a_i}:
 0\le d_i<1,\ \sum_{i=1}^m d_i=s\right\}.
 \tag{W.1}
\]
函数 $P_{\boldsymbol a}$ 连续、严格递减，且 $P_{\boldsymbol a}(0)=1$。因此当
\[
 P_{\boldsymbol a}(1)<\lambda<1
 \tag{W.2}
\]
时，存在唯一的 $s_{\boldsymbol a}(\lambda)\in(0,1)$ 满足
\[
 P_{\boldsymbol a}(s_{\boldsymbol a}(\lambda))=\lambda.
 \tag{W.3}
\]

\begin{lemma}[水位分配]\label{lem:water-filling}
若 $0<s<1$，存在唯一的 $\rho>0$ 使
\[
 \sum_{i=1}^m\left(1-\frac{a_i}{\rho}\right)_+=s.
 \tag{W.4}
\]
式 \textup{(W.1)} 的唯一最大点为
\[
 d_i^*(s)=\left(1-\frac{a_i}{\rho}\right)_+.
 \tag{W.5}
\]
\end{lemma}

\begin{proof}
在单纯形的相对内部，目标的对数为严格凹函数
\(\sum_i a_i\log(1-d_i)\)。KKT 条件给出：活跃坐标满足
\(a_i/(1-d_i)=\rho\)，非活跃坐标满足 $a_i\ge\rho$。于是得到
\textup{(W.4)}--\textup{(W.5)}。左边作为 $\rho$ 的函数连续递增，从 $0$ 取到大于 $1$ 的值，因而 $\rho$ 唯一。严格递减性和连续性也随之得到：增加总缺陷后，把一个正缺陷坐标减小即可严格增大乘积。
\end{proof}

设 $\mu$ 是非原子、局部有限的 Radon 测度，零质量方体一律略去，
\[
 \mathcal M_{c,\mu}^{\boldsymbol a}(f_1,\ldots,f_m)(x)
 =\sup_{r>0}\prod_{i=1}^m
 \left(\frac{1}{\mu(Q(x,r))}\int_{Q(x,r)}|f_i|\,d\mu\right)^{a_i}.
\]
定义共同中心常数
\[
 C_{c,\mu}^{\boldsymbol a}(\lambda)=
 \sup_{\substack{E_1,\ldots,E_m\ \mathrm{measurable}\\
 F=\cap_iE_i,\ 0<\mu(F)<\infty}}
 \frac{\mu\{\mathcal M_{c,\mu}^{\boldsymbol a}
 (\mathbf1_{E_1},\ldots,\mathbf1_{E_m})>\lambda\}}{\mu(F)}.
\]

\begin{theorem}[不等指数共同中心精确传递]\label{thm:weighted-transfer}
在 \textup{(W.2)} 的区间内，
\[
 C_{c,\mu}^{\boldsymbol a}(\lambda)=
 C_{c,\mu}^{\mathrm{int}}\bigl(1-s_{\boldsymbol a}(\lambda)\bigr).
 \tag{W.6}
\]
若 $d\mu=w\,dx$，上式适用于任意非零非负局部可积权。对外层模型，把左侧平均改为 Lebesgue 平均、光晕测量改为 $w\,dx$，并要求 $0<|F|<\infty$，仍有同一等式。
\end{theorem}

\begin{proof}
写 $s=s_{\boldsymbol a}(\lambda)$。若在共同见证方体 $Q$ 上
\[
 p_i=\frac{\mu(E_i\cap Q)}{\mu(Q)},\qquad d_i=1-p_i,
\]
且 $\prod_i p_i^{a_i}>\lambda$，若总缺陷至少为 $s$，可逐坐标减小缺陷到总和恰为 $s$；乘积只能增大，且至多为 $P_{\boldsymbol a}(s)=\lambda$，矛盾。因此定义 \textup{(W.1)} 给出
\(\sum_i d_i<s\)。又 $Q\setminus F\subset\bigcup_i(Q\setminus E_i)$，故
\[
 \frac{\mu(F\cap Q)}{\mu(Q)}
 \ge1-\sum_i d_i>1-s.
\]
这给出 \textup{(W.6)} 的“$\le$”方向。

反之，令 $K$ 为严格线性光晕
\(\{M_{c,\mu}\mathbf1_F>1-s\}$ 的紧子集。我们使用如下有限稳定事实：可选有限个参考方体以及正余量，使 $K$ 被相应的小中心邻域覆盖；每个邻域中的中心方体仍有 $F$-缺陷严格小于 $s$；并且参考方体及这些小中心扰动的对称差的 $\mu$-质量小到小于预先指定的余量。先把严格见证半径略微增大，避开至多可数个具有正质量边界的半径，使每个参考方体边界的 $\mu$-质量为零。严格性在充分小的增大下保持。此时这一事实由严格不等式、局部有限性、从上连续性和
\(\mu(Q(x,r)\mathbin\triangle Q(y,r))\to0\) 得到。

取包含全部微扰方体的有界盒 $D$。把有限参考方体所生成的 Boolean 原子上的 $D\setminus F$ 按下列比例切分：
\[
 \frac{d_1^*(s)}s,\ldots,\frac{d_m^*(s)}s.
\]
非原子性保证每个原子可作这种有限分割；在 $D$ 外把 $F^c$ 任意补入这些部分，得到 $F^c=A_1\mathbin{\dot\cup}\cdots\mathbin{\dot\cup}A_m$，并置 $E_i=\mathbb R^n\setminus A_i$。对未扰动的参考方体，
\[
 \frac{\mu(Q\setminus E_i)}{\mu(Q)}=
 \frac{d_i^*(s)}s\frac{\mu(Q\setminus F)}{\mu(Q)}<d_i^*(s)\quad(d_i^*(s)>0).
 \tag{W.7}
\]
有限稳定性中的余量使 \textup{(W.7)} 在所覆盖的全部中心扰动方体中仍严格成立。这里把 $D$ 取为包含所有微扰方体的一个有界盒；在该盒内按参考方体的 Boolean 原子分割。若 $d_i^*(s)=0$，令 $A_i\cap D=\varnothing$，故该坐标没有误差。于是每个相关方体满足
\[
 \prod_i\left(\frac{\mu(E_i\cap Q)}{\mu(Q)}\right)^{a_i}
 >\prod_i(1-d_i^*(s))^{a_i}=\lambda.
\]
因此 $K$ 包含于所构造的多线性光晕。由 Radon 内正则性令 $K$ 穷尽线性光晕，得到反向不等式。

外层版本完全相同：对 Lebesgue 测度作有限 Boolean 原子切分，再以 $w\,dx$ 的 Radon 内正则性穷尽光晕。
\end{proof}

\begin{corollary}[等指数特例]
当 $a_1=\cdots=a_m=1$ 时，
\[
 P_{\boldsymbol a}(s)=(1-s/m)^m,\qquad
 s_{\boldsymbol a}(\lambda)=m(1-\lambda^{1/m}),
\]
从而恢复等指数的共同中心缺陷传递。
\end{corollary}

\begin{remark}[外层临界发散]
若 $m\ge2$，则对每个非零局部可积权 $w$ 和
\(0<\lambda\le P_{\boldsymbol a}(1)\)，外层常数为无穷。证明取沿第一坐标的周期分割，其各部分的比例为 $d_i^*(1)$；把一个有正 $w$-质量的有限集 $F$ 放进某个正比例部分，并令 $E_i=(\mathbb R^n\setminus A_i)\cup F$。每个充分大的整数边长中心方体给出严格大于 $P_{\boldsymbol a}(1)$ 的共同见证，故光晕为全空间。有限总质量情形再令 $w(F)\downarrow0$。该论证改变内禀平均分母，不能推出内禀模型的临界发散。
\end{remark}
 纳入主稿；要求 amsmath, amssymb, amsthm。
\section{不等指数的共同中心传递}\label{sec:weighted-exponents}

本节把等指数乘积的缺陷坐标推广到任意正指数。该推广仍要求各因子在同一个中心方体上取平均；分别取上确界的多线性模型不适用本节结论。

令 $m\ge2$、$a_1,\ldots,a_m>0$，并在 $0\le s\le1$ 定义
\[
 P_{\boldsymbol a}(s)=
 \max\left\{\prod_{i=1}^m(1-d_i)^{a_i}:
 0\le d_i<1,\ \sum_{i=1}^m d_i=s\right\}.
 \tag{W.1}
\]
函数 $P_{\boldsymbol a}$ 连续、严格递减，且 $P_{\boldsymbol a}(0)=1$。因此当
\[
 P_{\boldsymbol a}(1)<\lambda<1
 \tag{W.2}
\]
时，存在唯一的 $s_{\boldsymbol a}(\lambda)\in(0,1)$ 满足
\[
 P_{\boldsymbol a}(s_{\boldsymbol a}(\lambda))=\lambda.
 \tag{W.3}
\]

\begin{lemma}[水位分配]\label{lem:water-filling}
若 $0<s<1$，存在唯一的 $\rho>0$ 使
\[
 \sum_{i=1}^m\left(1-\frac{a_i}{\rho}\right)_+=s.
 \tag{W.4}
\]
式 \textup{(W.1)} 的唯一最大点为
\[
 d_i^*(s)=\left(1-\frac{a_i}{\rho}\right)_+.
 \tag{W.5}
\]
\end{lemma}

\begin{proof}
在单纯形的相对内部，目标的对数为严格凹函数
\(\sum_i a_i\log(1-d_i)\)。KKT 条件给出：活跃坐标满足
\(a_i/(1-d_i)=\rho\)，非活跃坐标满足 $a_i\ge\rho$。于是得到
\textup{(W.4)}--\textup{(W.5)}。左边作为 $\rho$ 的函数连续递增，从 $0$ 取到大于 $1$ 的值，因而 $\rho$ 唯一。严格递减性和连续性也随之得到：增加总缺陷后，把一个正缺陷坐标减小即可严格增大乘积。
\end{proof}

设 $\mu$ 是非原子、局部有限的 Radon 测度，零质量方体一律略去，
\[
 \mathcal M_{c,\mu}^{\boldsymbol a}(f_1,\ldots,f_m)(x)
 =\sup_{r>0}\prod_{i=1}^m
 \left(\frac{1}{\mu(Q(x,r))}\int_{Q(x,r)}|f_i|\,d\mu\right)^{a_i}.
\]
定义共同中心常数
\[
 C_{c,\mu}^{\boldsymbol a}(\lambda)=
 \sup_{\substack{E_1,\ldots,E_m\ \mathrm{measurable}\\
 F=\cap_iE_i,\ 0<\mu(F)<\infty}}
 \frac{\mu\{\mathcal M_{c,\mu}^{\boldsymbol a}
 (\mathbf1_{E_1},\ldots,\mathbf1_{E_m})>\lambda\}}{\mu(F)}.
\]

\begin{theorem}[不等指数共同中心精确传递]\label{thm:weighted-transfer}
在 \textup{(W.2)} 的区间内，
\[
 C_{c,\mu}^{\boldsymbol a}(\lambda)=
 C_{c,\mu}^{\mathrm{int}}\bigl(1-s_{\boldsymbol a}(\lambda)\bigr).
 \tag{W.6}
\]
若 $d\mu=w\,dx$，上式适用于任意非零非负局部可积权。对外层模型，把左侧平均改为 Lebesgue 平均、光晕测量改为 $w\,dx$，并要求 $0<|F|<\infty$，仍有同一等式。
\end{theorem}

\begin{proof}
写 $s=s_{\boldsymbol a}(\lambda)$。若在共同见证方体 $Q$ 上
\[
 p_i=\frac{\mu(E_i\cap Q)}{\mu(Q)},\qquad d_i=1-p_i,
\]
且 $\prod_i p_i^{a_i}>\lambda$，若总缺陷至少为 $s$，可逐坐标减小缺陷到总和恰为 $s$；乘积只能增大，且至多为 $P_{\boldsymbol a}(s)=\lambda$，矛盾。因此定义 \textup{(W.1)} 给出
\(\sum_i d_i<s\)。又 $Q\setminus F\subset\bigcup_i(Q\setminus E_i)$，故
\[
 \frac{\mu(F\cap Q)}{\mu(Q)}
 \ge1-\sum_i d_i>1-s.
\]
这给出 \textup{(W.6)} 的“$\le$”方向。

反之，令 $K$ 为严格线性光晕
\(\{M_{c,\mu}\mathbf1_F>1-s\}$ 的紧子集。我们使用如下有限稳定事实：可选有限个参考方体以及正余量，使 $K$ 被相应的小中心邻域覆盖；每个邻域中的中心方体仍有 $F$-缺陷严格小于 $s$；并且参考方体及这些小中心扰动的对称差的 $\mu$-质量小到小于预先指定的余量。先把严格见证半径略微增大，避开至多可数个具有正质量边界的半径，使每个参考方体边界的 $\mu$-质量为零。严格性在充分小的增大下保持。此时这一事实由严格不等式、局部有限性、从上连续性和
\(\mu(Q(x,r)\mathbin\triangle Q(y,r))\to0\) 得到。

取包含全部微扰方体的有界盒 $D$。把有限参考方体所生成的 Boolean 原子上的 $D\setminus F$ 按下列比例切分：
\[
 \frac{d_1^*(s)}s,\ldots,\frac{d_m^*(s)}s.
\]
非原子性保证每个原子可作这种有限分割；在 $D$ 外把 $F^c$ 任意补入这些部分，得到 $F^c=A_1\mathbin{\dot\cup}\cdots\mathbin{\dot\cup}A_m$，并置 $E_i=\mathbb R^n\setminus A_i$。对未扰动的参考方体，
\[
 \frac{\mu(Q\setminus E_i)}{\mu(Q)}=
 \frac{d_i^*(s)}s\frac{\mu(Q\setminus F)}{\mu(Q)}<d_i^*(s)\quad(d_i^*(s)>0).
 \tag{W.7}
\]
有限稳定性中的余量使 \textup{(W.7)} 在所覆盖的全部中心扰动方体中仍严格成立。这里把 $D$ 取为包含所有微扰方体的一个有界盒；在该盒内按参考方体的 Boolean 原子分割。若 $d_i^*(s)=0$，令 $A_i\cap D=\varnothing$，故该坐标没有误差。于是每个相关方体满足
\[
 \prod_i\left(\frac{\mu(E_i\cap Q)}{\mu(Q)}\right)^{a_i}
 >\prod_i(1-d_i^*(s))^{a_i}=\lambda.
\]
因此 $K$ 包含于所构造的多线性光晕。由 Radon 内正则性令 $K$ 穷尽线性光晕，得到反向不等式。

外层版本完全相同：对 Lebesgue 测度作有限 Boolean 原子切分，再以 $w\,dx$ 的 Radon 内正则性穷尽光晕。
\end{proof}

\begin{corollary}[等指数特例]
当 $a_1=\cdots=a_m=1$ 时，
\[
 P_{\boldsymbol a}(s)=(1-s/m)^m,\qquad
 s_{\boldsymbol a}(\lambda)=m(1-\lambda^{1/m}),
\]
从而恢复等指数的共同中心缺陷传递。
\end{corollary}

\begin{remark}[外层临界发散]
若 $m\ge2$，则对每个非零局部可积权 $w$ 和
\(0<\lambda\le P_{\boldsymbol a}(1)\)，外层常数为无穷。证明取沿第一坐标的周期分割，其各部分的比例为 $d_i^*(1)$；把一个有正 $w$-质量的有限集 $F$ 放进某个正比例部分，并令 $E_i=(\mathbb R^n\setminus A_i)\cup F$。每个充分大的整数边长中心方体给出严格大于 $P_{\boldsymbol a}(1)$ 的共同见证，故光晕为全空间。有限总质量情形再令 $w(F)\downarrow0$。该论证改变内禀平均分母，不能推出内禀模型的临界发散。
\end{remark}
 之后纳入主稿；依赖 amsmath, amssymb, amsthm。
\begin{proposition}[有限总质量时的内禀临界发散]\label{prop:finite-mass-critical}
令 $m\ge2$、$a_1,\ldots,a_m>0$，并令 $\mu=w\,dx$ 是非零有限
Radon 测度（故 $\mu$ 非原子）。设 $P_{\boldsymbol a}$ 由
\textup{\rm(W.1)} 定义。对每个
\[
 0<\lambda\le P_{\boldsymbol a}(1)
\]
都有
\[
 C_{c,\mu}^{\boldsymbol a}(\lambda)=\infty.
\]
\end{proposition}

\begin{proof}
记 $M=\mu(\mathbb R^n)\in(0,\infty)$。在闭单纯形
$\{d_i\ge0:\sum_i d_i=1\}$ 上把目标连续延拓为
$\prod_i(1-d_i)^{a_i}$。取 $d_i=1/m$ 可知最大值为正；而边界上
若某个 $d_i=1$，该延拓值为零。因此存在最大点
$\boldsymbol d^*=(d_1^*,\ldots,d_m^*)$，满足
\[
 P_{\boldsymbol a}(1)=\prod_i(1-d_i^*)^{a_i}>0,
 \qquad 0\le d_i^*<1,\qquad \sum_i d_i^*=1.
\]
选取 $d_j^*>0$。非原子性给出 $\mathbb R^n=A_1\mathbin{\dot\cup}
\cdots\mathbin{\dot\cup}A_m$，其中 $\mu(A_i)=d_i^*M$（零比例部分
可取零测度）。任取 $0<t<d_j^*M$，再取 $F\subset A_j$ 使
$\mu(F)=t$，并置
\[
 E_i=A_i^c\cup F\quad(1\le i\le m).
\]
这时 $\bigcap_iE_i=F$，且
\[
 \prod_i\left(\frac{\mu(E_i)}M\right)^{a_i}
 =\left(1-d_j^*+\frac tM\right)^{a_j}
   \prod_{i\ne j}(1-d_i^*)^{a_i}
 >P_{\boldsymbol a}(1)\ge\lambda. \tag{FM.1}
\]
对每个 $x\in\mathbb R^n$，有 $Q(x,R)\uparrow\mathbb R^n$ 以及
$\mu(Q(x,R))\to M>0$。故左边的方体平均乘积收敛到 \textup{(FM.1)}
的左边；对充分大的（依赖于 $x$ 的）$R$，其严格大于 $\lambda$。
所以相应多线性光晕就是 $\mathbb R^n$，并且
\[
 \frac{\mu\{\mathcal M_{c,\mu}^{\boldsymbol a}
 (\mathbf1_{E_1},\ldots,\mathbf1_{E_m})>\lambda\}}{\mu(F)}
 =\frac Mt.
\]
令 $t\downarrow0$ 即得结论。
\end{proof}

\subsection{任意维乘积权的临界闭合}
设 $m\ge2$、$a_i>0$，并令 $w(x)=\prod_{\nu=1}^n w_\nu(x_\nu)$，其中每个因子非零、非负且局部可积。记 $d\mu_\nu=w_\nu\,dt$、$\mu=\otimes_\nu\mu_\nu$。算子及公共交集常数沿用前述定义，平均只取正有限质量的中心方体。
\begin{lemma}[累计质量坐标的整数周期恒等式]\label{lem:periodic}
设 $\mu_1(\R)=\infty$，并置
\[
 T(t)=\begin{cases}\mu_1((0,t]),&t\geq0,\\-\mu_1((t,0]),&t<0.\end{cases}
 \qquad J=T(\R).
\]
则 $T$ 连续、非减，$J$ 是无界区间，并且对任意 Borel 集
$B\subset J$ 有
\[
 \mu_1(T^{-1}(B))=|B|. \tag{P.2}
\]
给定 $d_1,\ldots,d_m\geq0$、$\sum_i d_i=1$，令 $S_i$ 是
$T^{-1}$ 下周期 $1$ 的半开分割之第 $i$ 块的逆像；即在每个
$[k,k+1)$ 中该块的长度为 $d_i$。若 $I=[u,v]$ 且
$T(v)-T(u)=N\in\mathbb N$，则
\[
 \mu_1(S_i\cap I)=d_iN=d_i\mu_1(I). \tag{P.3}
\]
\end{lemma}

\begin{proof}
局部绝对连续性给出 $T(v)-T(u)=\mu_1((u,v])$。先对半开区间
验证 (P.2)，再由单调类定理推广到 Borel 集；$T$ 的平坦区间的
逆像具有 $\mu_1$-质量零，故不影响此等式。由 $\mu_1(\R)=\infty$
知 $J$ 无界。周期分割在每一个完整长度 $1$ 的区间中给第 $i$
块长度 $d_i$；把 (P.2) 应用于 $[T(u),T(v)]\cap J$ 即得 (P.3)。
端点至多是零质量集合。
\end{proof}

\begin{theorem}[有无限因子的乘积权临界定理]\label{thm:product-critical}
设 $m\ge2$。若至少一个因子满足 $\mu_\nu(\R)=\infty$，则
\[
 C_{c,\mu}^{(1,\ldots,1)}(\lambda)=\infty
 \qquad\left(0<\lambda\leq\left(1-\frac1m\right)^m\right). \tag{P.4}
\]
更一般地，令
\[
 P_{\boldsymbol a}(1)=
 \max_{\substack{d_i\ge0\\\sum_i d_i=1}}
 \prod_{i=1}^m(1-d_i)^{a_i} \quad(a_i>0).
 \tag{P.5}
\]
则对 $0<\lambda\le P_{\boldsymbol a}(1)$ 有
\[
 C_{c,\mu}^{\boldsymbol a}(\lambda)=\infty. \tag{P.6}
\]
\end{theorem}

\begin{proof}
交换坐标后可设 $\mu_1(\R)=\infty$。取 (P.5) 的最大点
$\boldsymbol d=(d_1,\ldots,d_m)$。最大值为正（例如比较
$d_i=1/m$），故 $d_i<1$ 对所有 $i$ 成立；并选取一个 $j$ 使
$d_j>0$。引理的半开周期分割给出精确的集合分割
$\R=S_1\mathbin{\dot\cup}\cdots\mathbin{\dot\cup}S_m$，其中第 $i$ 块比例为
$d_i$。定义
\[
 A_i=S_i\times\R^{n-1}.
\]
每个 $w_\nu$ 非零，故 $\mu_\nu(\R)>0$。$J$ 无界又说明
$\mu_1(S_j)=\infty$；从而 $A_j$ 含一个有界 Borel 集 $F$，满足
\[
 0<f:=\mu(F)<\infty. \tag{P.7}
\]
令 $E_i=A_i^c\cup F$。由于 $F\subset A_j$，有
$\cap_iE_i=F$。

固定 $x\in\R^n$。函数
\[
 g_x(r)=\mu_1([x_1-r,x_1+r])
=T(x_1+r)-T(x_1-r)
\]
连续、非减且趋于无穷。先选 $r_0$，使 $Q(x,r_0)$ 包含 $F$ 且其余
每个坐标区间已有正质量；再取整数 $N>g_x(r_0)$，由介值性选到
$r\ge r_0$ 使 $g_x(r)=N$。于是 $Q(x,r)$ 包含 $F$ 且其每个坐标
区间都有正的 $\mu_\nu$-质量。后一个条件可行，因为每个
$\mu_\nu(\R)>0$。局部可积性给出这些质量有限。由 (P.3) 和
乘积结构，精确地有
\[
 \frac{\mu(A_i\cap Q(x,r))}{\mu(Q(x,r))}=d_i. \tag{P.8}
\]
于是
\[
 \frac{\mu(E_i\cap Q(x,r))}{\mu(Q(x,r))}
 =\begin{cases}
 1-d_i,&i\ne j,\\
 1-d_j+f/\mu(Q(x,r)),&i=j.
 \end{cases}
\]
所有右边的底数均正，故
\[
 \prod_{i=1}^m\left(\frac{\mu(E_i\cap Q(x,r))}{\mu(Q(x,r))}\right)^{a_i}
 >\prod_{i=1}^m(1-d_i)^{a_i}=P_{\boldsymbol a}(1)\ge\lambda. \tag{P.9}
\]
这对每一个 $x$ 成立，故相应光晕为全空间。又
$\mu(\R^n)=\prod_\nu\mu_\nu(\R)=\infty$，以 (P.7) 归一化即得
(P.6)。等指数时，AM--GM 给 (P.5) 的唯一最大点
$d_i=1/m$，并给出 (P.4)。
\end{proof}

\begin{remark}[适用范围]
本证明只使用一维累计质量坐标、周期的完整整数块和乘积测度分解；
它并不处理一般非乘积的无穷质量权。若所有因子总质量均有限，则
本证明中 ``全空间光晕有无穷质量'' 的最后一步失效；该情形应改由
有限非原子测度的全局等分构造单独处理。
\end{remark}


\begin{corollary}[全部乘积权]
对于任意维非零乘积权及任意正指数，在 $0<\lambda\le P_{\boldsymbol a}(1)$ 时，内禀共同中心多集合常数均为无穷。
\end{corollary}
\begin{proof}
若至少一个因子质量无穷，用本节定理；若所有因子总质量有限，则乘积测度总质量有限，使用命题 \ref{prop:finite-mass-critical}。特别地，取 $n=1$ 得任意一维权的不等指数版本。
\end{proof}

\section{应用与开放问题}
由等指数传递和逐权一阶项，固定 $m$ 与 $w$ 时
\[
 C^{(m),\rm int}_{c,n,w}(1-\epsilon)=1+2^n\epsilon+o_{m,w}(\epsilon).
\]
这是因为 $s_m(1-\epsilon)=\epsilon+O_m(\epsilon^2)$。此外，若外层集中模量满足 $\Omega_w(t)\le A_0t^\sigma$，代入 $t=s_m(\lambda)$ 即得到多集合外层估计；不等指数时改代 $s_{\boldsymbol a}(\lambda)$。
若一个中心几何族对于所有局部有限测度的弱型常数统一不超过 $B$，上面的测度扰动证明还给出 $C_\mu(a)\le1+B(1-a)/a$。此推广只使用弱型性质，不要求把 Besicovitch 交叠性质当作覆盖选择定理。


外层幂权的匹配尖锐首项、固定权二阶极限、高维非乘积的一般无限总质量权的临界以下内禀多线性发散，以及 restricted strong type 的完整尖锐范数仍是开放问题；势论或 PDE 应用还须独立验证共同中心方体见证接口。
\enlargethispage{2\baselineskip}
\begin{thebibliography}{9}
\bibitem{Aldaz2020} J. M. Aldaz, \emph{Kissing numbers and the centered maximal operator}, arXiv:1811.10478v2, 2020.
\bibitem{HagelsteinParissis2016} P. A. Hagelstein and I. Parissis, \emph{Weighted Solyanik estimates for the Hardy--Littlewood maximal operator and embedding of $A_\infty$ into $A_p$}, J. Geom. Anal. 26 (2016), 924--946.
\end{thebibliography}\end{document}
